% 2.1 — brief en documento/investigacion/2-1-sistemas-erp.md
\subsection{SISTEMAS ERP}
\subsubsection{DEFINICIÓN Y CARACTERÍSTICAS}

Los sistemas de planificación de recursos empresariales (\textit{enterprise resource planning}, ERP) no cuentan con una definición única. \textcite{Klaus2000} encontraron visiones divergentes sobre el fenómeno y concluyeron que era poco probable alcanzar una definición ampliamente consensuada. Existe, sin embargo, un núcleo común. \textcite{Davenport1998} los describe como paquetes de software comerciales que prometen la integración fluida de toda la información que circula por una empresa: la financiera y contable, la de recursos humanos, la de la cadena de suministro y la de los clientes. \textcite{Shi2003} lo formulan en términos operativos: un entorno integrado que permite articular las principales funciones de gestión sobre una única base de datos común, de modo que la información se comparta entre ellas.

Ese núcleo se explica por su origen. El ERP nació en la industria manufacturera \parencite{Shi2003}; \textcite{Jacobs2007} reconstruyen esa historia y el papel de los principales proveedores. \textcite{Klaus2000} sitúan el origen del término en la planificación de requerimientos de materiales (MRP) y la planificación de recursos de manufactura (MRP II): cuando las empresas advirtieron que la rentabilidad y la satisfacción del cliente eran objetivos de toda la organización, la integración se extendió a finanzas, ventas, distribución y recursos humanos, y a esa solución integral se la llamó ERP\@.

De esa caracterización se desprenden tres rasgos. El primero es la integración: una sola base de datos elimina la duplicación y la inconsistencia entre sistemas heredados dispersos, el atractivo que \textcite{Davenport1998} identifica frente a sistemas incompatibles. El segundo es la organización en módulos funcionales que comparten esa base y pueden adoptarse de forma progresiva. El tercero es la estandarización de procesos: el paquete trae incorporada una forma de operar y, como advierte el mismo autor, impone su propia lógica sobre la estrategia, la organización y la cultura de la empresa, por lo que su adopción debe ser responsabilidad de la gerencia general y no solo de los tecnólogos.

Los beneficios y los riesgos están documentados. \textcite{Shang2002}, a partir de 233 casos y entrevistas con gerentes de 34 organizaciones, consolidan los beneficios en cinco dimensiones (operacional, gerencial, estratégica, de infraestructura de TI y organizacional) y los evalúan en perspectiva longitudinal, años después de la implementación. Del lado de los riesgos, \textcite{Umble2003} describen al ERP como un sistema de información altamente complejo, cuya implementación es difícil y costosa, y señalan que muchas implementaciones han sido calificadas como fracasos por no alcanzar los objetivos corporativos predeterminados.

Frente al paquete comercial existe la alternativa del desarrollo propio. \textcite{Olsen2007} sostienen que para muchas pequeñas y medianas empresas el software desarrollado en casa es una opción óptima: ofrece control total, flexibilidad y adaptación a requisitos específicos, con costos a menudo comparables a los de un ERP estándar. Cuando el software soporta el núcleo del negocio, el desarrollo propio otorga ventaja competitiva, mientras que las estructuras rígidas de los paquetes pueden amenazar la dinámica organizacional. Es la situación de ISI Mustang: sus procesos centrales (centros de costo, certificaciones por avance, requisiciones, viáticos) se apoyan en un ERP desarrollado a medida, que hoy se reconstruye sobre NestJS y PostgreSQL a partir de un portal anterior en PHP y MySQL, y que este trabajo toma a la vez como objeto de mejora y como fuente de datos.

\subsubsection{ERP EN EMPRESAS DE INGENIERÍA}

La aplicación del ERP a las empresas de ingeniería y construcción exige revisar el supuesto manufacturero. \textcite{Shi2003} advierten que la naturaleza particular de la industria de la construcción impide adoptar directamente los sistemas orientados a manufactura y proponen las características fundamentales de un ERP para construcción: funciones de gestión de proyectos, técnicas de planificación avanzada y estandarización de procesos. En estas empresas la unidad de gestión deja de ser el producto repetitivo y pasa a ser el proyecto, con alcance, presupuesto, plazo y cliente propios, y el sistema debe seguir cada uno por separado sin perder la visión consolidada de la empresa.

Aun con esas adaptaciones, la adopción en el sector no es sencilla. \textcite{Chung2008} reconocen que los ERP ofrecen muchos beneficios a la industria de ingeniería y construcción, pero que muchas firmas dudan en adoptarlos por su alto costo, sus incertidumbres y sus riesgos; en respuesta identifican factores críticos de éxito y fracaso e indicadores para evaluar el éxito del ERP, y proponen un modelo que los relaciona. El trabajo traslada al sector la discusión genérica de \textcite{Umble2003}, y su enfoque asume que el resultado de una implementación depende de factores organizacionales y no solo técnicos.

Una segunda línea conecta el ERP con el control de desempeño de los proyectos. \textcite{Skibniewski2009} se preguntan qué áreas de proceso de la construcción no pueden ser cubiertas por el ERP a la hora de recolectar indicadores clave de desempeño (KPI), y lo responden con una encuesta a los responsables de ERP de grandes firmas de ingeniería y construcción de Estados Unidos. El planteamiento del trabajo supone que una parte de los procesos de proyecto queda fuera del alcance de un ERP estándar, lo que refuerza el argumento de \textcite{Olsen2007} a favor del desarrollo propio cuando el núcleo del negocio es el proyecto.

Lo anterior se corresponde con la forma en que ISI Mustang ha organizado su propio sistema. En el ERP de la empresa la unidad que articula toda la operación es el centro de costo, que puede representar un proyecto, un servicio o un área de estructura y tiene código, responsable, moneda y vigencia propios. Sobre esa unidad se registran las horas del personal, las requisiciones y las órdenes de compra emitidas a los proveedores, los viáticos y sus rendiciones, y el presupuesto, con una línea base separada del seguimiento. El ingreso ejecutado de cada proyecto se registra mediante certificaciones, cuyo flujo distingue la solicitud por parte del gerente del centro de costo, el registro y el cobro.

El portal anterior, cuyos procesos se analizan en el Capítulo III, calculaba además indicadores de seguimiento y control por proyecto, entre ellos el desvío de costo, el desvío de ingresos, el avance del proyecto y el avance de trabajo, a partir del monto de la orden de compra emitida por el cliente, lo ejecutado y el flujo de caja planificado. Esos indicadores son el antecedente directo de las desviaciones que el módulo predictivo busca anticipar, y los datos con los que se calculaban son los que se exploran en el Capítulo III\@.

\subsubsection{ERP EN EL SECTOR DE PETRÓLEO, GAS Y MINERÍA}

La literatura sobre ERP en petróleo, gas y minería es más escasa que la de manufactura o construcción y está compuesta sobre todo por estudios de caso, cercanos a los sectores en los que ISI Mustang ejecuta sus proyectos. En petróleo y gas, \textcite{AlJafari2018} analizan, con un enfoque mixto, cómo la implementación del ERP influye en la eficiencia y la productividad de las empresas del sector en el Sultanato de Omán. Parten de que el papel del ERP en esa industria no puede negarse y concluyen recomendando su implementación como palanca de crecimiento.

Los retos organizacionales aparecen con más nitidez en \textcite{Ali2023}, que mediante investigación-acción estudian las barreras culturales a la implementación de ERP en el sector de petróleo y gas de Medio Oriente, desde la planificación hasta la posimplementación. Identifican como causas del retraso el conflicto entre equipos, la autoridad gerencial, la falta de una cultura de TI y de compromiso con la capacitación, y la tecnofobia. \textcite{Gool2018}, por su parte, plantean el dilema de cuándo personalizar el ERP y cuándo cambiar los procesos de negocio para ajustarlos a la funcionalidad estándar. En una petrolera multinacional africana encuentran que la personalización continúa después de la implementación y describen prácticas para gestionarla: capacitación, eliminación sistemática de modificaciones y procesos de aprobación.

Para la minería, \textcite{Larasati2023} encuestan a la gerencia y el personal de una empresa minera indonesia y ordenan los factores críticos de fracaso del ERP con el método TOPSIS (\textit{Technique for Order of Preference by Similarity to Ideal Solution}); los dos más importantes resultan ser la mala comprensión de los procesos de negocio de la organización y una reingeniería de procesos deficiente. El hallazgo coincide con lo observado en petróleo y gas: en las industrias extractivas el ERP fracasa o se retrasa menos por la tecnología que por el conocimiento de los procesos, la cultura de uso y la gestión continua de la adaptación. Cabe suponer que lo mismo vale para las empresas EPC que trabajan para ellas, aunque ninguno de los casos revisados las estudia directamente.

Tres lecciones de esa literatura orientan este proyecto: documentar los flujos de negocio antes de sistematizarlos, criterio que se sigue en el análisis de procesos del portal anterior; asumir la personalización como un proceso continuo, lo que es coherente con la metodología iterativa e incremental de la construcción; y anticipar que la resistencia al uso señalada por \textcite{Ali2023} puede afectar también a un módulo predictivo, de ahí la validación con usuarios clave prevista en el Capítulo IV\@.

Queda, por último, el papel del ERP que más interesa a este trabajo: el de fuente de datos. Un sistema que durante más de diez años ha registrado horas, compras, presupuestos e ingresos facturados por proyecto, complementado con las planillas de control que la empresa lleva fuera del sistema, acumula un histórico del que pueden aprenderse patrones. \textcite{MLdrivenERP2023} revisan la integración del aprendizaje automático con los sistemas ERP: sus algoritmos extraen patrones complejos de grandes conjuntos de datos y permiten pronósticos más precisos, con un énfasis creciente en la interpretabilidad de los resultados. Ese es el enlace con las secciones que siguen: el aprendizaje automático aporta los métodos, la analítica predictiva el marco de aplicación, y el ERP de ISI Mustang, tanto el histórico del portal anterior en MySQL como los proyectos activos del ERP nuevo en PostgreSQL, los datos.
